% Datasets of a snapshot file and the fields declared from them (invented example).
\begin{tikzpicture}[fig]
\tikzset{ds/.style={font=\scriptsize\ttfamily, anchor=west, inner sep=1.5pt},
         fl/.style={draw, rounded corners, font=\scriptsize, align=left, inner sep=3pt, text width=3.6cm, anchor=west}}
\draw[rounded corners, cGrid, thick] (-0.2,-3.05) rectangle (3.3,0.55);
\node[lab, anchor=west] at (-0.1,0.3) {snapshot file (HDF5)};
\node[ds] (C)  at (0,-0.2) {C\quad (N, N)};
\node[ds] (V)  at (0,-0.75) {V\quad (2, N, N)};
\node[ds] (px) at (0,-1.3) {px\quad (N, N)};
\node[ds] (py) at (0,-1.75) {py\quad (N, N)};
\node[ds] (Qx) at (0,-2.3) {Qxx\quad (N, N)};
\node[ds] (Qy) at (0,-2.75) {Qxy\quad (N, N)};
\node[lab, anchor=west] at (5.0,0.3) {fields declared};
\node[fl, fill=cP2!10] (frho) at (5.0,-0.2) {\textbf{rho}: scalar, relative};
\node[fl, fill=cP0!10] (fv) at (5.0,-0.75) {\textbf{v}: vector (velocity)};
\node[fl, fill=cP1!10] (fp) at (5.0,-1.52) {\textbf{p}: vector (polarity),\\ components px, py};
\node[fl, fill=black!6] (fQ) at (5.0,-2.52) {\textbf{Q}: nematic,\\ components Qxx, Qxy};
\draw[->] (C) -- (frho);
\draw[->] (V) -- (fv);
\draw[->] (px.east) -- ++(0.4,0) |- (fp.west);
\draw[->] (py.east) -- ++(0.4,0) |- (fp.west);
\draw[->] (Qx.east) -- ++(0.4,0) |- (fQ.west);
\draw[->] (Qy.east) -- ++(0.4,0) |- (fQ.west);
\node[small, text width=4.2cm, align=left, anchor=north west] at (9.0,0.35)
  {\textbf{Why grouping matters.} A rotation by $90^\circ$ sends $(p_x,p_y)\to(-p_y,p_x)$: $p_x$ alone
   becomes another field, so its descriptors are not invariant, and the augmentations of the
   network would be wrong. A vector is one field with two components.\\[3pt]
   Today: one dataset per field (\code{V}, \code{P}, \code{Q}); assembling \code{px}, \code{py}
   is planned.};
\end{tikzpicture}
